\documentclass[11pt]{article}

\usepackage[utf8]{inputenc}
\usepackage[T1]{fontenc}
\usepackage{lmodern}
\usepackage[margin=1in]{geometry}
\usepackage{amsmath,amssymb}
\usepackage{booktabs}
\usepackage{tabularx}
\usepackage{array}
\usepackage{enumitem}
\usepackage{hyperref}

\hypersetup{hidelinks}
\setlength{\parindent}{0pt}
\setlength{\parskip}{0.6em}
\renewcommand{\arraystretch}{1.2}

\title{Real-Life Applications of Continuous Probability Density Functions\\with Simple Geometric Shapes}
\author{}
\date{}

\begin{document}
	
	\maketitle
	
	\section{Introduction}
	
	Continuous probability density functions are often introduced through simple geometric shapes such as rectangles, triangles, and piecewise line segments. These models are useful because probabilities can be found as areas under the density curve, often using only geometry rather than integration.
	
	This report presents five real-world examples in which a continuous random variable can be reasonably modeled by one of the following shapes:
	
	\begin{itemize}[leftmargin=*]
		\item Uniform (rectangle)
		\item Triangular
		\item Piecewise-constant
		\item Piecewise-linear
	\end{itemize}
	
	For each example, the report gives the real-life context, the random variable and its units, a reasonable domain, a proposed normalized density function, one probability question in context, and a full solution using area formulas only. Each example is supported by a credible source that makes the modeling choice realistic, even if the proposed density is still a simplified approximation.
	
	\section{Geometric idea behind the calculations}
	
	If $X$ is a continuous random variable with density $f(x)$, then:
	
	\[
	f(x)\ge 0
	\qquad\text{and}\qquad
	\text{total area under }f(x)=1.
	\]
	
	Probabilities are areas:
	
	\[
	P(a\le X\le b)=\text{area under }f(x)\text{ from }a\text{ to }b.
	\]
	
	In this report, every probability is found using only these area formulas:
	
	\begin{itemize}[leftmargin=*]
		\item Rectangle: $A=\text{base}\times \text{height}$
		\item Triangle: $A=\frac12(\text{base})(\text{height})$
		\item Trapezoid: $A=(\text{width})\dfrac{\text{left height}+\text{right height}}{2}$
	\end{itemize}
	
	\section{Five real-world examples}
	
	\subsection{Example 1: Uniform density --- digital instrument resolution error}
	
	\textbf{Context.} A digital measuring instrument displays values with a fixed resolution. If the instrument rounds to the nearest $0.1^\circ\text{C}$, then the true value can lie anywhere within half a step above or below the displayed value. A standard metrology reference models this type of uncertainty with a rectangular distribution.
	
	\textbf{Random variable and units.} Let $E$ be the rounding or resolution error, in degrees Celsius.
	
	\textbf{Reasonable domain.}
	\[
	E\in[-0.05,\,0.05].
	\]
	
	\textbf{Best-fitting density shape and why.} A uniform density is appropriate because, within one display interval, there is no reason to prefer one point over another.
	
	\textbf{Proposed density function.}
	\[
	f_E(e)=
	\begin{cases}
		10, & -0.05\le e\le 0.05,\\
		0, & \text{otherwise}.
	\end{cases}
	\]
	
	This is valid because the total area is
	\[
	(0.1)(10)=1.
	\]
	
	\textbf{Probability question in context.} What is the probability that the error is at most $0.02^\circ\text{C}$ in magnitude?
	
	\[
	P(|E|\le 0.02).
	\]
	
	\textbf{Solution using area only.} The favorable interval is $[-0.02,0.02]$, so the base is
	\[
	0.02-(-0.02)=0.04.
	\]
	
	The density height on this interval is $10$, so the probability is the area of a rectangle:
	\[
	P(|E|\le 0.02)=0.04\times 10=0.40.
	\]
	
	\textbf{Interpretation.} There is a $40\%$ chance that the rounding error falls within $\pm 0.02^\circ\text{C}$.
	
	\textbf{Limitation.} This model assumes all values inside the display interval are equally likely. In practice, real measurement processes may not be perfectly uniform within each step.
	
	\textbf{Supporting source.} Reference \,[1].
	
	\subsection{Example 2: Triangular density --- duration estimate in project planning}
	
	\textbf{Context.} In project planning and risk analysis, a task is often described by three values: an optimistic time, a most likely time, and a pessimistic time. This naturally leads to a triangular model.
	
	\textbf{Random variable and units.} Let $T$ be the duration of a task, in hours.
	
	\textbf{Reasonable domain.} Suppose:
	\[
	\text{minimum}=2,\qquad \text{most likely}=5,\qquad \text{maximum}=8.
	\]
	
	Then
	\[
	T\in[2,8].
	\]
	
	\textbf{Best-fitting density shape and why.} A triangular density fits because the middle value is the most plausible, while the minimum and maximum are still possible but less likely.
	
	\textbf{Proposed density function.}
	\[
	f_T(t)=
	\begin{cases}
		\dfrac{t-2}{9}, & 2\le t\le 5,\\[6pt]
		\dfrac{8-t}{9}, & 5<t\le 8,\\
		0, & \text{otherwise}.
	\end{cases}
	\]
	
	This is valid because the graph is one triangle with base $6$ and height
	\[
	f_T(5)=\frac{5-2}{9}=\frac13,
	\]
	so the total area is
	\[
	\frac12(6)\left(\frac13\right)=1.
	\]
	
	\textbf{Probability question in context.} What is the probability that the task lasts between $4$ and $6$ hours?
	
	\[
	P(4\le T\le 6).
	\]
	
	\textbf{Solution using area only.} It is easiest to use the complement:
	\[
	P(4\le T\le 6)=1-P(T<4)-P(T>6).
	\]
	
	The left tail from $2$ to $4$ is a triangle with base $2$ and height
	\[
	f_T(4)=\frac{4-2}{9}=\frac29.
	\]
	
	So
	\[
	P(T<4)=\frac12(2)\left(\frac29\right)=\frac29.
	\]
	
	By symmetry around $t=5$,
	\[
	P(T>6)=\frac29.
	\]
	
	Therefore,
	\[
	P(4\le T\le 6)=1-\frac29-\frac29
	=1-\frac49
	=\frac59.
	\]
	
	\textbf{Interpretation.} The probability that the task takes between $4$ and $6$ hours is
	\[
	\frac59\approx 0.556.
	\]
	
	\textbf{Limitation.} The triangular model is simple, but real task durations may have longer tails or more complicated shapes than one straight-sided peak.
	
	\textbf{Supporting source.} Reference \,[2].
	
	\subsection{Example 3: Piecewise-constant density --- waiting time when transit frequency changes by time window}
	
	\textbf{Context.} Public transit service frequency can change during the day. In one published transit plan, a trolley route runs every $20$ minutes most of the time, but during a certain window an additional loop makes service effectively every $10$ minutes at most locations. This creates different waiting-time regimes.
	
	\textbf{Random variable and units.} Let $W$ be the waiting time for the next trolley, in minutes.
	
	\textbf{Reasonable domain.}
	\[
	W\in[0,20].
	\]
	
	\textbf{Modeling idea.} If service is regular and a rider arrives at a random time during a window with headway $H$, then the waiting time in that window is reasonably modeled as uniform on $[0,H]$.
	
	Using the published schedule:
	\begin{itemize}[leftmargin=*]
		\item effective $10$-minute service for $300$ minutes of the day,
		\item effective $20$-minute service for $530$ minutes of the day,
		\item total operating window $830$ minutes.
	\end{itemize}
	
	So the time-window weights are
	\[
	P(\text{10-minute window})=\frac{300}{830}=\frac{30}{83},
	\qquad
	P(\text{20-minute window})=\frac{530}{830}=\frac{53}{83}.
	\]
	
	\textbf{Best-fitting density shape and why.} This gives a mixture of two uniform models, so the overall density is constant on one interval and then drops to a lower constant level on the next interval. That is a piecewise-constant density.
	
	\textbf{Proposed density function.}
	\[
	f_W(w)=
	\begin{cases}
		\frac{30}{83}\cdot \frac{1}{10}+\frac{53}{83}\cdot \frac{1}{20}, & 0\le w\le 10,\\[8pt]
		\frac{53}{83}\cdot \frac{1}{20}, & 10<w\le 20,\\
		0, & \text{otherwise}.
	\end{cases}
	\]
	
	Simplifying,
	\[
	f_W(w)=
	\begin{cases}
		\dfrac{113}{1660}, & 0\le w\le 10,\\[6pt]
		\dfrac{53}{1660}, & 10<w\le 20,\\
		0, & \text{otherwise}.
	\end{cases}
	\]
	
	This is normalized because
	\[
	10\left(\frac{113}{1660}\right)+10\left(\frac{53}{1660}\right)
	=\frac{1130+530}{1660}=1.
	\]
	
	\textbf{Probability question in context.} What is the probability that a randomly arriving rider waits at most $5$ minutes?
	
	\[
	P(W\le 5).
	\]
	
	\textbf{Solution using area only.} From $0$ to $5$, the density is constant at $\dfrac{113}{1660}$, so the probability is the area of a rectangle:
	\[
	P(W\le 5)=5\left(\frac{113}{1660}\right)=\frac{565}{1660}=\frac{113}{332}.
	\]
	
	\textbf{Interpretation.}
	\[
	P(W\le 5)=\frac{113}{332}\approx 0.340.
	\]
	
	So the chance of waiting no more than $5$ minutes is about $34.0\%$.
	
	\textbf{Limitation.} This assumes perfectly regular service and random rider arrival. Real transit systems often have delays, bunching, and uneven spacing.
	
	\textbf{Supporting source.} Reference \,[3].
	
	\subsection{Example 4: Piecewise-linear density --- total error from two independent uniform errors}
	
	\textbf{Context.} In measurement, more than one source of resolution error can affect a final reported value. If two independent error components are each modeled as uniform, then their sum produces a trapezoidal density, which is piecewise linear.
	
	\textbf{Random variable and units.} Let $Z$ be the total measurement error, in grams, where
	\[
	E_1\sim \text{Uniform}[-0.5,0.5],\qquad
	E_2\sim \text{Uniform}[-1.5,1.5],
	\]
	and
	\[
	Z=E_1+E_2.
	\]
	
	\textbf{Reasonable domain.}
	\[
	Z\in[-2,2].
	\]
	
	\textbf{Best-fitting density shape and why.} The sum of two uniform errors with different widths produces a trapezoidal density. A trapezoid is piecewise linear because it rises linearly, stays flat, and then falls linearly.
	
	\textbf{Proposed density function.}
	\[
	f_Z(z)=
	\begin{cases}
		\dfrac{z+2}{3}, & -2\le z\le -1,\\[6pt]
		\dfrac{1}{3}, & -1<z<1,\\[6pt]
		\dfrac{2-z}{3}, & 1\le z\le 2,\\
		0, & \text{otherwise}.
	\end{cases}
	\]
	
	This is normalized by area:
	\begin{itemize}[leftmargin=*]
		\item left triangle: $\dfrac12(1)\left(\dfrac13\right)=\dfrac16$,
		\item middle rectangle: $(2)\left(\dfrac13\right)=\dfrac23$,
		\item right triangle: $\dfrac16$.
	\end{itemize}
	
	Total area:
	\[
	\frac16+\frac23+\frac16=1.
	\]
	
	\textbf{Probability question in context.} What is the probability that the total error is at most $0.5$ grams in magnitude?
	
	\[
	P(|Z|\le 0.5).
	\]
	
	\textbf{Solution using area only.} The interval $[-0.5,0.5]$ lies completely inside the flat middle part, where the height is $\dfrac13$. So the probability is the area of a rectangle with base
	\[
	0.5-(-0.5)=1
	\]
	and height $\dfrac13$:
	\[
	P(|Z|\le 0.5)=1\left(\frac13\right)=\frac13.
	\]
	
	\textbf{Interpretation.}
	\[
	P(|Z|\le 0.5)=\frac13\approx 0.333.
	\]
	
	\textbf{Limitation.} This model assumes the two component errors are independent and purely due to resolution. Real measurement error may also include bias, drift, or calibration effects.
	
	\textbf{Supporting source.} Reference \,[1].
	
	\subsection{Example 5: Piecewise-linear density --- time of day of peak rainfall in a valley}
	
	\textbf{Context.} Studies of rainfall in the Medell\'in Valley report a diurnal cycle in which rainfall tends to peak during afternoon hours in part of the year. A simple piecewise-linear density can approximate a gradual rise toward the afternoon peak and a gradual decline afterward.
	
	\textbf{Random variable and units.} Let $H$ be the local hour of day at which the daily rainfall peak occurs, measured in hours from $0$ to $24$.
	
	\textbf{Reasonable domain.}
	\[
	H\in[0,24].
	\]
	
	\textbf{Best-fitting density shape and why.} A piecewise-linear density is suitable because the chance of the peak does not jump suddenly from hour to hour. Instead, it can rise and fall gradually.
	
	\textbf{Proposed density function.} Define
	\[
	a=\frac{1}{50},
	\qquad
	b=\frac{31}{250}.
	\]
	
	Then let
	\[
	f_H(h)=
	\begin{cases}
		\dfrac{1}{50}, & 0\le h\le 10,\\[6pt]
		\dfrac{1}{50}+\dfrac{13}{750}(h-10), & 10<h\le 16,\\[8pt]
		\dfrac{31}{250}-\dfrac{13}{500}(h-16), & 16<h\le 20,\\[8pt]
		\dfrac{1}{50}, & 20<h\le 24,\\
		0, & \text{otherwise}.
	\end{cases}
	\]
	
	\textbf{Normalization check by area only.} The total area is the sum of two rectangles and two trapezoids:
	
	\[
	\text{rectangles: } (10+4)\left(\frac{1}{50}\right)=14\left(\frac{1}{50}\right)=\frac{14}{50},
	\]
	
	\[
	\text{trapezoid }[10,16]: 6\cdot \frac{a+b}{2}=3(a+b),
	\]
	
	\[
	\text{trapezoid }[16,20]: 4\cdot \frac{a+b}{2}=2(a+b).
	\]
	
	So total area is
	\[
	14a+3(a+b)+2(a+b)=19a+5b.
	\]
	
	Substituting $a=\dfrac{1}{50}$ and $b=\dfrac{31}{250}$:
	\[
	19\left(\frac{1}{50}\right)+5\left(\frac{31}{250}\right)
	=\frac{19}{50}+\frac{31}{50}=1.
	\]
	
	\textbf{Probability question in context.} What is the probability that the rainfall peak occurs between 15:00 and 18:00?
	
	\[
	P(15\le H\le 18).
	\]
	
	\textbf{Solution using area only.} Split the interval into $[15,16]$ and $[16,18]$.
	
	For $h=15$,
	\[
	f_H(15)=\frac{1}{50}+\frac{13}{750}(5)
	=\frac{1}{50}+\frac{65}{750}
	=\frac{8}{75}.
	\]
	
	For $h=16$,
	\[
	f_H(16)=\frac{31}{250}.
	\]
	
	So the area from $15$ to $16$ is a trapezoid of width $1$:
	\[
	A_{15,16}=1\cdot \frac{\frac{8}{75}+\frac{31}{250}}{2}.
	\]
	
	Using denominator $750$,
	\[
	\frac{8}{75}=\frac{80}{750},\qquad
	\frac{31}{250}=\frac{93}{750},
	\]
	so
	\[
	A_{15,16}=\frac12\left(\frac{173}{750}\right)=\frac{173}{1500}.
	\]
	
	Now find the height at $h=18$:
	\[
	f_H(18)=\frac{31}{250}-\frac{13}{500}(2)
	=\frac{31}{250}-\frac{13}{250}
	=\frac{18}{250}
	=\frac{9}{125}.
	\]
	
	The area from $16$ to $18$ is a trapezoid of width $2$:
	\[
	A_{16,18}=2\cdot \frac{\frac{31}{250}+\frac{9}{125}}{2}
	=\frac{31}{250}+\frac{9}{125}
	=\frac{31}{250}+\frac{18}{250}
	=\frac{49}{250}.
	\]
	
	Therefore,
	\[
	P(15\le H\le 18)=A_{15,16}+A_{16,18}
	=\frac{173}{1500}+\frac{49}{250}.
	\]
	
	Since
	\[
	\frac{49}{250}=\frac{294}{1500},
	\]
	we get
	\[
	P(15\le H\le 18)=\frac{173}{1500}+\frac{294}{1500}
	=\frac{467}{1500}.
	\]
	
	\textbf{Interpretation.}
	\[
	P(15\le H\le 18)=\frac{467}{1500}\approx 0.311.
	\]
	
	So the model gives about a $31.1\%$ chance that the daily peak happens between 3 p.m.\ and 6 p.m.
	
	\textbf{Limitation.} The real rainfall pattern may be more complex than one smooth afternoon peak. This density is a simplified teaching model built to reflect the broad observed trend rather than every meteorological detail.
	
	\textbf{Supporting source.} Reference \,[4].
	
	\section{Which models are most realistic and which are most simplified?}
	
	The most simplified models are the uniform and triangular densities. They are very useful when only limited information is available:
	
	\begin{itemize}[leftmargin=*]
		\item The uniform model is appropriate when only the lower and upper bounds matter.
		\item The triangular model is appropriate when a minimum, a most likely value, and a maximum are known.
	\end{itemize}
	
	The more realistic models in this set are the piecewise-constant and piecewise-linear densities because they can reflect changing conditions or gradual trends:
	
	\begin{itemize}[leftmargin=*]
		\item The piecewise-constant waiting-time model reflects two different service regimes in the same day.
		\item The piecewise-linear models can capture gradual increases and decreases, which is often more realistic than sudden jumps.
	\end{itemize}
	
	Even so, none of these models should be treated as exact descriptions of reality. They are useful approximations chosen because they are mathematically simple and easy to interpret geometrically.
	
	\section{Overall limitations and assumptions}
	
	Across all five examples, the main assumptions are:
	
	\begin{itemize}[leftmargin=*]
		\item The variable is treated as continuous, even if real measurements may be rounded or recorded in discrete steps.
		\item The chosen density is an approximation rather than a fully data-fitted model.
		\item External sources support the realism of the context, but the exact numerical density in each example is still a simplified classroom model.
		\item The probability questions are solved with geometry only, so each example is intentionally designed to use simple shapes.
	\end{itemize}
	
	\section{Summary table}
	
	\small
	\begin{tabularx}{\textwidth}{>{\raggedright\arraybackslash}p{2.7cm} >{\raggedright\arraybackslash}p{2.2cm} >{\raggedright\arraybackslash}p{2.2cm} >{\raggedright\arraybackslash}p{2.3cm} X}
		\toprule
		\textbf{Context} & \textbf{Variable} & \textbf{Domain} & \textbf{Density shape} & \textbf{Why it fits} \\
		\midrule
		Digital instrument resolution error & $E$ in $^\circ\text{C}$ & $[-0.05,0.05]$ & Uniform & All values inside one display interval are treated as equally plausible. \\
		\addlinespace
		Project task duration with optimistic, most likely, and pessimistic times & $T$ in hours & $[2,8]$ & Triangular & The middle value is most plausible, while the endpoints are less likely but still possible. \\
		\addlinespace
		Transit waiting time with different service windows & $W$ in minutes & $[0,20]$ & Piecewise-constant & Different time windows create different constant waiting-time densities. \\
		\addlinespace
		Total error from two independent uniform errors & $Z$ in grams & $[-2,2]$ & Piecewise-linear (trapezoidal) & Adding two uniform error sources gives a trapezoidal density with linear sides. \\
		\addlinespace
		Hour of daily rainfall peak & $H$ in hours & $[0,24]$ & Piecewise-linear & The probability rises and falls gradually through the day rather than changing abruptly. \\
		\bottomrule
	\end{tabularx}
	\normalsize
	
	\section{Conclusion}
	
	Simple density shapes are valuable because they connect probability to geometry in a very direct way. Uniform and triangular models are especially good for quick approximations when limited information is available. Piecewise-constant and piecewise-linear models are more flexible and often better at reflecting real changes in conditions over time.
	
	The key lesson is that a probability density function does not need to be complicated to be useful. If it is normalized, nonnegative, and reasonable for the context, it can provide meaningful answers and clear interpretations. For students, these geometric models are a strong bridge between real-world contexts and the formal study of continuous probability.
	
	\section*{References}
	
	\begin{enumerate}[leftmargin=*]
		\item Joint Committee for Guides in Metrology. \emph{Evaluation of Measurement Data --- Guide to the Expression of Uncertainty in Measurement (JCGM 100:2008)}. Available at \url{https://www.bipm.org/documents/20126/2071204/JCGM_100_2008_E.pdf}.
		
		\item NASA. \emph{Appendix G: Cost Risk and Uncertainty Methodologies}. Available at \url{https://www.nasa.gov/wp-content/uploads/2020/11/ceh_appg.pdf}.
		
		\item Wave Transit. \emph{Wave Short Range Transit Plan}. Available at \url{https://www.wavetransit.com/wp-content/uploads/2016/08/SRTP_2012.pdf}.
		
		\item Bedoya-Soto, J. M., Poveda, G., and Vieco, D. \emph{Seasonal Shift of the Diurnal Cycle of Rainfall Over Medellin's Valley, Central Andes of Colombia (1998--2005)}. \emph{Frontiers in Earth Science}. Available at \url{https://www.frontiersin.org/journals/earth-science/articles/10.3389/feart.2019.00092/full}.
	\end{enumerate}
	
\end{document}